\chapter{微分算子层\index{ziceng@子层|hyperpage}\index{weifensuanziceng@微分算子层|hyperpage}}
\label{chap:01.1}

\section{扭微分算子层\index{niuweifensuanziceng@扭微分算子层|hyperpage}}
\label{sec:01.1.1}

设 $X$ 是特征 $0$ 的代数闭域 $k$ 上的光滑代数簇，$\calO_X$ 是 $X$ 的结构层，$T_X$ 是 $X$ 的切层，$\D_X$ 是 $X$ 上的微分算子层。考虑二元组 $(A,i)$ 构成的范畴，其中 $A$ 是 $X$ 上带单位元的结合 $k$-代数层，$i:\calO_X\to A$ 是带单位元的 $k$-代数态射。层 $\D_X$ 连同自然包含 $i_X:\calO_X\to\D_X$ 是该范畴的一个对象。若二元组 $(\D,i)$ 局部同构于 $(\D_X,i_X)$，即若 $X$ 有开覆盖 $\{U\}$，使得对其中每个 $U$ 都有
$$(\D|_U,i|_U)\cong(\D_U,i_U),$$
则称 $(\D,i)$ 是 $X$ 上的\emph{扭微分算子层}（twisted sheaf of differential operators）。

下面讨论 $X$ 上扭微分算子层的自然参数化。先作一些准备。

\begin{lemma}\label{lem:01.1.1}
设 $\varphi$ 是 $(\D_X,i_X)$ 的自同态。则存在 $X$ 上的闭 $1$-形式 $\omega$，使得对任意局部向量场\index{jubuxiangliangchang@局部向量场|hyperpage} $\xi\in T_X$ 有
$$\varphi(\xi)=\xi-\omega(\xi),$$
并且 $\varphi$ 由 $\omega$ 完全确定。特别地，$\varphi$ 是 $(\D_X,i_X)$ 的自同构。
\end{lemma}

\begin{proof}
设 $f\in\calO_X$，$\xi\in T_X$。则
$$[\varphi(\xi),f]=[\varphi(\xi),\varphi(f)]=\varphi([\xi,f])=\varphi(\xi(f))=\xi(f).$$
将上式作用于函数 $1$，得
$$\varphi(\xi)(f)=\xi(f)+f\,\varphi(\xi)(1).$$
因此可令 $\omega(\xi)=-\varphi(\xi)(1)$。显然 $\omega$ 是 $X$ 上的 $1$-形式。又对 $\xi,\eta\in T_X$ 有
\begin{align*}
\omega([\xi,\eta])&=-\varphi([\xi,\eta])(1)=-(\varphi(\xi)\varphi(\eta)-\varphi(\eta)\varphi(\xi))(1)\\
&=\varphi(\xi)(\omega(\eta))-\varphi(\eta)(\omega(\xi))=\xi(\omega(\eta))-\eta(\omega(\xi)).
\end{align*}
于是
$$d\omega(\xi\wedge\eta)=\xi(\omega(\eta))-\eta(\omega(\xi))-\omega([\xi,\eta])=0,$$
即 $d\omega=0$，$\omega$ 是闭形式。$\varphi$ 由 $\omega$ 完全确定，这一点是显然的。此外，$\varphi$ 保持 $\D_X$ 的滤过\index{lvguo@滤过|hyperpage}，且诱导的分次自同态 $\Gr\varphi:\Gr\D_X\to\Gr\D_X$ 是恒等态射。由此立见 $\varphi$ 是自同构。
\end{proof}

由 \ref{lem:01.1.1}，$(\D_X,i_X)$ 的每个自同构 $\varphi$ 都决定 $X$ 上的一个闭 $1$-形式 $\omega$。显然，这个对应是乘法群 $\Aut(\D_X,i_X)$ 到闭 $1$-形式加法群 $Z^1(X)$ 的单同态。

\begin{lemma}\label{lem:01.1.2}
自然态射 $\Aut(\D_X,i_X)\to Z^1(X)$ 是同构。
\end{lemma}

\begin{proof}
只需再证该态射是满的。设 $\omega$ 是 $X$ 上的闭 $1$-形式。定义映射 $\varphi:T_X\to\D_X$ 为
$$\varphi(\xi)=\xi-\omega(\xi),\qquad \xi\in T_X.$$
显然 $\varphi$ 满足 [9] 中的条件 (ii) 与 (iii)。又因 $\omega$ 闭，对 $\xi,\eta\in T_X$ 有
\begin{align*}
\varphi([\xi,\eta])&=[\xi,\eta]-\omega([\xi,\eta])=[\xi,\eta]-\xi(\omega(\eta))+\eta(\omega(\xi))\\
&=[\xi-\omega(\xi),\,\eta-\omega(\eta)]=[\varphi(\xi),\varphi(\eta)],
\end{align*}
即条件 (i) 也满足。因此 $\varphi$ 可延拓为 $(\D_X,i_X)$ 的自同态。由 \ref{lem:01.1.1}，$\varphi$ 实际上是自同构。
\end{proof}

设 $(\D,i)$ 是 $X$ 上的扭微分算子层。则存在开覆盖 $\calU=(U_j;1\le j\le n)$，使得对 $1\le j\le n$ 都有 $(\D,i)|_{U_j}\cong(\D_{U_j},i_{U_j})$。对每个 $j$ 取定同构
$$\psi_j:(\D,i)|_{U_j}\longrightarrow(\D_{U_j},i_{U_j}).$$
于是存在 $(\D_{U_j\cap U_k},i_{U_j\cap U_k})$ 的自同构 $\varphi_{jk}$，使图表
\begin{center}
\begin{tikzcd}
(\D,i)|_{U_j\cap U_k} \arrow[rr,"\psi_j"] \arrow[dr,"\psi_k"'] & & (\D_{U_j\cap U_k},i_{U_j\cap U_k}) \arrow[dl,"\varphi_{jk}"]\\
 & (\D_{U_j\cap U_k},i_{U_j\cap U_k}) &
\end{tikzcd}
\end{center}
交换，即 $\psi_j=\varphi_{jk}\circ\psi_k$。由 \ref{lem:01.1.2}，存在 $U_j\cap U_k$ 上的闭 $1$-形式 $\omega_{jk}$，它决定 $\varphi_{jk}$。若 $U_j\cap U_k\cap U_l\neq\emptyset$，则在其上有
$$\psi_j=\varphi_{jk}\circ\psi_k=\varphi_{jk}\circ\varphi_{kl}\circ\psi_l,$$
故在 $U_j\cap U_k\cap U_l$ 上 $\varphi_{jl}=\varphi_{jk}\circ\varphi_{kl}$。
\begin{annote}
底本此处把 $\varphi_{jk}\circ\varphi_{kl}$ 印作 $\psi_{jk}\circ\psi_{kl}$。$\psi_{jk}$ 在全书没有定义，且上一行已给出 $\varphi_{jl}=\varphi_{jk}\circ\varphi_{kl}$，故按上下文改正。[来源：底本 p.2--3，上下文自证]
\end{annote}
这又推出对 $\xi\in T_{U_j\cap U_k\cap U_l}$ 有
$$\varphi_{jl}(\xi)=\xi-\omega_{jl}(\xi)=(\varphi_{jk}\circ\varphi_{kl})(\xi)=\varphi_{jk}(\xi-\omega_{kl}(\xi))=\xi-\omega_{jk}(\xi)-\omega_{kl}(\xi),$$
即
$$\omega_{jl}=\omega_{jk}+\omega_{kl}\qquad\text{于 } U_j\cap U_k\cap U_l.$$

设 $Z^1_X$ 是 $X$ 上闭 $1$-形式构成的层，$C^\bullet(\calU,Z^1_X)$ 是 $Z^1_X$ 关于覆盖 $\calU$ 的 \v{C}ech 复形。则 $\omega=(\omega_{jk};1\le j<k\le n)$ 是 $C^1(\calU,Z^1_X)$ 的元素，且 $d\omega=0$，即 $\omega\in Z^1(\calU,Z^1_X)$。现设另取一组局部同构
$$\psi'_j:(\D,i)|_{U_j}\longrightarrow(\D_{U_j},i_{U_j}),\qquad 1\le j\le n.$$
这将给出另一组 $(\varphi'_{jk};1\le j<k\le n)$ 与另一个 $\omega'\in Z^1(\calU,Z^1_X)$。再用一次 \ref{lem:01.1.2}，可得 $(\D_{U_j},i_{U_j})$ 的自同构 $\sigma_j$（$1\le j\le n$）以及与之相伴的闭 $1$-形式 $\rho_j$（$1\le j\le n$），使得
$$\psi'_j=\sigma_j\circ\psi_j,\qquad 1\le j\le n.$$
显然 $\rho=(\rho_j;1\le j\le n)$ 是 $C^0(\calU,Z^1_X)$ 的元素。现在在 $U_j\cap U_k$ 上有
$$\sigma_j\circ\varphi_{jk}\circ\psi_k=\sigma_j\circ\psi_j=\psi'_j=\varphi'_{jk}\circ\psi'_k=\varphi'_{jk}\circ\sigma_k\circ\psi_k,$$
故 $\sigma_j\circ\varphi_{jk}=\varphi'_{jk}\circ\sigma_k$。这给出
$$\rho_j+\omega_{jk}=\omega'_{jk}+\rho_k$$
于 $U_j\cap U_k$。由此可见 $\omega'-\omega=d\rho$。因此扭微分算子层 $(\D,i)$ 决定了 $H^1(\calU,Z^1_X)$ 中的一个元素。

于是得到良定义的映射
$$t:(\D,i)\longmapsto t(\D,i)\in H^1(X,Z^1_X),$$
它把扭微分算子层的同构类映到 $X$ 的以 $Z^1_X$ 为系数的第一 \v{C}ech 上同调群。

\begin{theorem}\label{thm:01.1.3}
映射 $t$ 给出 $X$ 上扭微分算子层的同构类与 $H^1(X,Z^1_X)$ 的元素之间的双射。
\end{theorem}

\begin{proof}
先证 $t$ 是单射。设 $\D$ 与 $\D'$ 是两个扭微分算子层，且 $t(\D)=t(\D')$。则二者决定同一个开覆盖 $\calU=(U_j;1\le j\le n)$ 与 $\omega,\omega'\in Z^1(\calU,Z^1_X)$，它们在 $H^1(X,Z^1_X)$ 中定义同一个元素；同时还决定局部同构族
$$\psi_j:\D|_{U_j}\to\D_{U_j},\qquad \varphi'_j:\D'|_{U_j}\to\D_{U_j}\qquad (1\le j\le n),$$
如前面的讨论所述。必要时用 $\calU$ 的加细，可设 $\omega-\omega'=d\rho$，其中 $\rho=(\rho_j;1\le j\le n)\in C^0(\calU,Z^1_X)$。设 $\sigma_j:\D_{U_j}\to\D_{U_j}$ 是由 $\rho_j$ 决定的自同构（$1\le j\le n$）。则 $\varphi''_j=\sigma_j\circ\varphi'_j$（$1\le j\le n$）是一族局部同构 $\varphi''_j:\D'|_{U_j}\to\D_{U_j}$，并具有性质
$$\sigma_j\circ\varphi'_{jk}\circ\psi'_k=\sigma_j\circ\psi'_j=\psi''_j=\varphi''_{jk}\circ\psi''_k=\varphi''_{jk}\circ\sigma_k\circ\psi'_k,$$
从而 $\sigma_j\circ\varphi'_{jk}=\varphi''_{jk}\circ\sigma_k$，即在 $U_j\cap U_k$ 上
$$\varphi''_{jk}=\sigma_j\circ\varphi'_{jk}\circ\sigma_k^{-1}.$$
这意味着在 $U_j\cap U_k$ 上
$$\omega''_{jk}=\omega'_{jk}+\rho_j-\rho_k,$$
即 $\omega''=\omega'+d\rho=\omega$。最终得到在 $U_j\cap U_k$ 上 $\varphi_{jk}=\varphi''_{jk}$。定义局部同构
$$\theta_j=\psi''^{-1}_j\circ\psi_j:\D|_{U_j}\longrightarrow\D'|_{U_j},\qquad 1\le j\le n.$$
则在 $U_j\cap U_k$ 上有
$$\theta_j=\psi''^{-1}_j\circ\psi_j=(\varphi_{jk}\circ\psi''_k)^{-1}\circ\varphi_{jk}\circ\psi_k=\psi''^{-1}_k\circ\psi_k=\theta_k,$$
故 $\theta$ 可延拓为 $\D$ 到 $\D'$ 的整体同构。

满性的证明是标准的\emph{粘接\index{zhanjie@粘接|hyperpage}}（recollement）论证，其中用到 \ref{lem:01.1.2}。
\end{proof}

下面描述 $X$ 上一类扭微分算子层的构造。设 $L$ 是 $X$ 上的可逆 $\calO_X$-模，$\D_L$ 是 $L$ 的全体微分自同态构成的层。由于 $L$ 局部同构于 $\calO_X$，$\D_L$ 是 $X$ 上的扭微分算子层。设 $\calO_X^*$ 是 $\calO_X$ 中可逆元构成的子层。众所周知，Picard 群 $\Pic(X)$ 等于 $H^1(X,\calO_X^*)$。对数导数\index{duishudaoshu@对数导数|hyperpage}给出交换群层的同态
$$d\log:\calO_X^*\longrightarrow Z^1_X,\qquad d\log f=f^{-1}df,\quad f\in\calO_X^*,$$
它诱导上同调群的态射 $H^p(d\log):H^p(X,\calO_X^*)\to H^p(X,Z^1_X)$。设 $i(L)$ 是 $L$ 在 $H^1(X,\calO_X^*)$ 中对应的元素。则有下述结果。

\begin{proposition}\label{prop:01.1.4}
对 $X$ 上任意可逆 $\calO_X$-模 $L$，
$$t(\D_L)=H^1(d\log)(i(L)).$$
\end{proposition}

\begin{proof}
取 $X$ 的开覆盖 $\calU=(U_i;1\le i\le n)$，使得对一切 $1\le i\le n$ 都有 $L|_{U_i}\cong\calO_X|_{U_i}$。记相应的 $\calO_X$-模同构为
$$\alpha_i:L|_{U_i}\longrightarrow\calO_X|_{U_i},\qquad 1\le i\le n,$$
并由
$$s_{jk}=\alpha_j(\alpha_k^{-1}(1)),\qquad 1\le j<k\le n,$$
定义 $\calO_X^*$ 在 $U_j\cap U_k$ 上的截面 $s_{jk}$。于是对 $L|_{U_j\cap U_k}$ 的截面 $s$ 有
$$\alpha_j(s)=\alpha_j(\alpha_k^{-1}(\alpha_k(s)))=s_{jk}\,\alpha_k(s),$$
即 $s=(s_{jk};1\le j<k\le n)$ 是代表 $i(L)$ 的 $1$-上闭链\index{shangbilian@上闭链|hyperpage}。此外，$\alpha_i$ 通过
$$\psi_i(u)=\alpha_i\circ u\circ\alpha_i^{-1},\qquad 1\le i\le n,$$
定义了 $\D_L|_{U_i}$ 到 $\D_X|_{U_i}$ 的同构 $\psi_i$。这意味着对任意 $f\in\calO_X|_{U_j\cap U_k}$ 有
\begin{align*}
\psi_j(u)(f)&=(\alpha_j\circ u\circ\alpha_j^{-1})(f)=\alpha_j\bigl(u(\alpha_j^{-1}(f))\bigr)\\
&=s_{jk}\,\alpha_k\bigl(u(\alpha_k^{-1}(s_{jk}^{-1}f))\bigr)=s_{jk}\,\psi_k(u)(s_{jk}^{-1}f).
\end{align*}
因此 $\psi_j(u)=s_{jk}\,\psi_k(u)\,s_{jk}^{-1}$，从而对任意 $D\in\D_X|_{U_j\cap U_k}$ 有
$$\varphi_{jk}(D)=s_{jk}\,D\,s_{jk}^{-1}.$$
设 $\zeta\in T_X|_{U_j\cap U_k}$，则
$$\varphi_{jk}(\zeta)=s_{jk}\,\zeta\,s_{jk}^{-1}=\zeta-s_{jk}^{-1}\zeta(s_{jk}),$$
即 $\omega_{jk}=s_{jk}^{-1}ds_{jk}$。因此代表 $t(\D_L)$ 的 $1$-上闭链 $\omega$ 由 $(s_{jk}^{-1}ds_{jk};1\le j<k\le n)$ 给出。
\end{proof}

下面研究函子性。设 $X,Y$ 是两个光滑代数簇，$\phi:X\to Y$ 是代数簇的态射，$\D$ 是 $Y$ 上的扭微分算子层。则 $\D_{X\to Y}=\phi^*(\D)$ 关于左乘是 $\calO_X$-模，关于右乘是右 $\phi^{-1}(\D)$-模。记 $\D_\phi$ 为右 $\phi^{-1}(\D)$-模 $\D_{X\to Y}$ 的全体微分自同态构成的层。显然，$\D_\phi$ 是 $X$ 上的结合代数层。此外还有自然的代数层态射 $i_\phi:\calO_X\to\D_\phi$。于是由 [9] 可知，二元组 $(\D_\phi,i_\phi)$ 局部同构于 $(\D_X,i_X)$，即 $\D_\phi$ 是 $X$ 上的扭微分算子层。由 \ref{thm:01.1.3}，$\D_\phi$ 决定 $H^1(X,Z^1_X)$ 中的元素 $t(\D_\phi)$。现在来计算 $t(\D_\phi)$。

先需要一个提升\index{tisheng@提升|hyperpage}引理。设 $(C,i)$ 与 $(\D,j)$ 是 $Y$ 上的两个扭微分算子层，$\psi:(C,i)\to(\D,j)$ 是同构。于是 $\psi$ 关于左乘与右乘所给的结构都是 $\calO_Y$-模同构。因此 $\psi$ 诱导 $\calO_X$-模的同构
$$\phi^*(\psi):C_{X\to Y}\longrightarrow\D_{X\to Y}.$$
又若 $u\in C_{X\to Y}$、$v\in\phi^{-1}(C)$，则
$$\phi^*(\psi)(uv)=\phi^*(\psi)(u)\,\phi^{-1}(\psi)(v).$$
由此可见，对任意 $z\in C_\phi$、$u\in\D_{X\to Y}$、$v\in\phi^{-1}(\D)$ 有
\begin{align*}
(\phi^*(\psi)z\phi^*(\psi^{-1}))(uv)&=\phi^*(\psi)\bigl(z(\phi^*(\psi^{-1})(u)\,\phi^{-1}(\psi^{-1})(v))\bigr)\\
&=\phi^*(\psi)\bigl((z\phi^*(\psi^{-1})(u))\,\phi^{-1}(\psi^{-1})(v)\bigr)\\
&=(\phi^*(\psi)z\phi^*(\psi^{-1}))(u)\,v,
\end{align*}
即 $\phi^*(\psi)z\phi^*(\psi^{-1})\in\D_\phi$。于是若置
$$\phi^\#(\psi)(z)=\phi^*(\psi)\,z\,\phi^*(\psi^{-1}),$$
则 $\phi^\#(\psi):C_\phi\to\D_\phi$ 是 $X$ 上 $k$-代数层的同构。显然 $j_\phi=\phi^\#(\psi)\circ i_\phi$。因此 $\phi^\#(\psi)$ 是扭微分算子层的同构，称之为 $\psi$ 的\emph{提升}（lifting）。此外，对任意 $z\in C_\phi$、$u\in C_{X\to Y}$ 有
$$(\phi^\#(\psi)(z))\,\phi^*(\psi)(u)=\phi^*(\psi)(zu).$$

现在考虑 $\D_Y$ 的自同构 $\alpha$ 这一特殊情形。由 \ref{lem:01.1.2}，它由 $Y$ 上的闭 $1$-形式 $\omega$ 决定。由 [9]，存在 $(\D_X,i_X)$ 与 $(\D_Y^\phi,i_{Y,\phi})$ 之间的自然同构 $\delta$。我们想计算 $\alpha$ 的提升 $\phi^\#(\alpha)$ 在 $\D_X$ 上诱导的自同构；更确切地说，计算由 \ref{lem:01.1.2} 它所决定的 $X$ 上的闭 $1$-形式。

设 $x\in X$，$U$ 是 $\phi(x)\in Y$ 的小开邻域，使得可以找到 $f_i\in\calO_Y(U)$（$1\le i\le\dim Y$），使 $df_i$（$1\le i\le\dim Y$）构成自由 $\calO_U$-模 $T^*Y|_U$ 的一组基。设 $\partial_i$ 是 $T_Y|_U$ 中的对偶基。则由 [9] 可见，对 $x$ 附近的局部向量场 $\xi$ 有
\begin{align*}
\gamma(\phi^\#(\alpha)(\xi))(1)&=\beta\bigl(\phi^\#(\alpha)(\xi)(1\otimes 1)\bigr)=\beta\bigl(\phi^*(\alpha)(\xi(1\otimes 1))\bigr)\\
&=\beta\Bigl(\phi^*(\alpha)\Bigl(\sum_i \xi(f_i\circ\phi)\otimes\phi^{-1}\partial_i\Bigr)\Bigr)=\beta\Bigl(\sum_i \xi(f_i\circ\phi)\otimes\phi^{-1}(\alpha\partial_i)\Bigr)\\
&=-\sum_i \xi(f_i\circ\phi)\,\omega(\partial_i)=-\omega\Bigl(\sum_i \xi(f_i\circ\phi)\,\partial_i\Bigr)=-(\phi^*\omega)(\xi),
\end{align*}
这里我们用 $\phi^*\omega$ 记 $\omega$ 在 $X$ 上诱导的 $1$-形式。因此 $\phi^*\omega$ 是与 $\phi^\#(\alpha)$ 相伴的闭 $1$-形式。

现在把这一事实用于 $t(\D_\phi)$ 的计算。设 $\calU=(U_i;1\le i\le n)$ 是 $Y$ 的开覆盖，
\begin{annote}
底本此处称 $\calU$ 是「$X$ 的开覆盖」，但下标 $U_i$ 随后要用于 $\D_Y$ 的局部平凡化，而且下一句定义 $V_i=\phi^{-1}(U_i)$ 为 $X$ 的开覆盖，故 $\calU$ 应是 $Y$ 的开覆盖。[来源：底本 p.5，上下文自证]
\end{annote}
$\psi_i:(\D,i)|_{U_i}\to(\D_Y,i_Y)|_{U_i}$ 是相应的同构。如前，对 $1\le j<k\le n$ 记 $\varphi_{jk}$ 为 $(\D_Y,i_Y)|_{U_j\cap U_k}$ 的自同构，使得 $\psi_j=\varphi_{jk}\circ\psi_k$。设 $V=(V_i;1\le i\le n)$ 是由 $V_i=\phi^{-1}(U_i)$（$1\le i\le n$）给出的 $X$ 的开覆盖。则提升 $\phi^\#(\psi_i)$ 是 $(\D_\phi,i_\phi)|_{V_i}$ 到 $(\D_X,i_X)|_{V_i}$ 的同构，而提升 $\phi^\#(\varphi_{jk})$ 是 $(\D_Y,i_Y)|_{U_j\cap U_k}$ 的自同构，满足
$$\phi^\#(\psi_j)=\phi^\#(\varphi_{jk})\circ\phi^\#(\psi_k),\qquad 1\le j<k\le n.$$
\begin{annote}
底本此处把两个 $\phi^\#$ 都印作 $\phi^*$。按上文提升的定义，$\psi_j$ 与 $\psi_k$ 的像都是扭微分算子层（而不是 $\D_{X\to Y}$）的同构，故应为 $\phi^\#$。[来源：底本 p.5，按定义校正]
\end{annote}
由前一段的讨论可见：若 $\omega=(\omega_{jk};1\le j<k\le n)$ 是 $Y$ 上与 $\D$ 相应的闭 $1$-形式上闭链，则 $(\phi^*\omega_{jk};1\le j<k\le n)$ 是 $X$ 上与 $\D_\phi$ 相应的闭 $1$-形式上闭链。映射 $\omega\mapsto\phi^*\omega$ 把 $Y$ 上的闭 $1$-形式映为 $X$ 上的闭 $1$-形式，从而诱导向量空间层的态射 $\phi^{-1}(Z^1_Y)\to Z^1_X$。由 [6, 3.2.2]，这个态射对每个 $p\in\Z_+$ 诱导线性映射
$$Z^p(\phi):H^p(Y,Z^1_Y)\longrightarrow H^p(X,Z^1_X).$$
因此前面的讨论实际上证明了下面的结果。

\begin{proposition}\label{prop:01.1.5}
设 $\phi:X\to Y$ 是光滑代数簇的态射，$\D$ 是 $Y$ 上的扭微分算子层。则
$$t(\D_\phi)=Z^1(\phi)(t(\D)).$$
\end{proposition}

此外，这一构造与态射的复合相容。

\begin{proposition}\label{prop:01.1.6}
设 $\phi:X\to Y$ 与 $\psi:Y\to Z$ 是光滑代数簇的态射，$\D$ 是 $Z$ 上的扭微分算子层。则 $\D_{\psi\circ\phi}$ 自然同构于 $(\D_\psi)_\phi$。
\end{proposition}

\begin{proof}
显然
\begin{align*}
(\psi\circ\phi)^*(\D)&=\phi^*(\psi^*(\D))=\phi^*(\D_{Y\to Z})=\calO_X\otimes_{\phi^{-1}(\calO_Y)}\phi^{-1}(\D_{Y\to Z})\\
&=\calO_X\otimes_{\phi^{-1}(\calO_Y)}\phi^{-1}(\D_\psi)\otimes_{\phi^{-1}(\D_\psi)}\phi^{-1}(\D_{Y\to Z})\\
&=(\D_\psi)_{X\to Y}\otimes_{\phi^{-1}(\D_\psi)}\phi^{-1}(\D_{Y\to Z}),
\end{align*}
作为 $\calO_X$-模与右 $(\psi\circ\phi)^{-1}(\D)$-模。又 $(\D_\psi)_\phi$ 在末式第一个因子上的作用显然与 $(\psi\circ\phi)^{-1}(\D)$ 的右作用交换。因此存在从 $(\D_\psi)_\phi$ 到 $\D_{\psi\circ\phi}$ 的自然态射。由 \ref{lem:01.1.1}，这个态射是扭微分算子层的同构。
\end{proof}

以下我们通过这个同构把 $(\D_\psi)_\phi$ 与 $\D_{\psi\circ\phi}$ 等同起来。

另一个要讨论的构造是：用可逆 $\calO_X$-模 $L$ 去扭光滑代数簇 $X$ 上的扭微分算子层 $\D$。若把 $\D$ 作为关于左乘的 $\calO_X$-模，可作 $\calO_X$-模 $L\otimes_{\calO_X}\D$。层 $L\otimes_{\calO_X}\D$ 关于第二个因子上的右乘是右 $\D$-模。因此可以考虑右 $\D$-模 $L\otimes_{\calO_X}\D$ 的局部微分自同态构成的层 $\D_L$。它显然是 $X$ 上的 $k$-代数层。又因 $L\otimes_{\calO_X}\D$ 是 $\calO_X$-模，有自然同态 $i_L:\calO_X\to\D_L$。我们断言 $(\D_L,i_L)$ 是 $X$ 上的扭微分算子层。取 $X$ 的开覆盖 $\calU=(U_i;1\le i\le n)$，使 $L|_{U_i}\cong\calO_X|_{U_i}$ 且 $\D|_{U_i}\cong\D_X|_{U_i}$。于是作为 $\calO_X|_{U_i}$-模，$(L\otimes_{\calO_X}\D)|_{U_i}$ 同构于 $\D_X|_{U_i}$。又在这个同构下，$(L\otimes_{\calO_X}\D)|_{U_i}$ 上的右 $\D|_{U_i}$-作用对应于 $\D_X|_{U_i}$ 上的右 $\D_X|_{U_i}$-作用。这诱导微分自同态层的同构，从而把 $\D_L|_{U_i}$ 与「把 $\D_X|_{U_i}$ 看作右 $\D_X|_{U_i}$-模」所得的微分自同态层等同起来。显然，这个代数层自然同构于 $\D_X|_{U_i}$。因此 $\D_L$ 是 $X$ 上的扭微分算子层。称之为 $\D$ 关于 $L$ 的\emph{扭}（twist）。

我们从下述结果开始研究扭。

\begin{lemma}\label{lem:01.1.7}
设 $L$ 是光滑代数簇 $X$ 上的可逆 $\calO_X$-模。则微分算子层 $\D_X$ 的扭 $(\D_X)_L$ 自然同构于 $\D_L$。
\end{lemma}

\begin{proof}
设 $I_X$ 是 $\D_X$ 中由 $T_X$ 生成的左理想。则有 $\calO_X$-模的正合列\index{zhenghelie@正合列|hyperpage}
$$0\longrightarrow I_X\longrightarrow\D_X\longrightarrow\calO_X\longrightarrow 0,$$
与 $L$ 作张量积得
$$0\longrightarrow L\otimes_{\calO_X}I_X\longrightarrow L\otimes_{\calO_X}\D_X\longrightarrow L\longrightarrow 0.$$
由 $(\D_X)_L$ 的构造可知这是 $(\D_X)_L$-模的正合列。因此存在从 $(\D_X)_L$ 到 $\D_L$ 的自然态射。由 \ref{lem:01.1.1}，它是扭微分算子层的同构。
\end{proof}

以下我们通过这个同构把 $(\D_X)_L$ 与 $\D_L$ 等同起来。

\begin{proposition}\label{prop:01.1.8}
设 $\D$ 是光滑代数簇 $X$ 上的扭微分算子层，$L$ 是可逆 $\calO_X$-模。则
$$t(\D_L)=t(\D)+H^1(d\log)(i(L)).$$
\end{proposition}

\begin{proof}
取 $X$ 的开覆盖 $\calU=(U_i;1\le i\le n)$ 以及局部同构
$$\alpha_i:L|_{U_i}\to\calO_X|_{U_i},\qquad \psi_i:(\D,i)|_{U_i}\to(\D_X,i_X)|_{U_i},\qquad 1\le i\le n.$$
如同 \ref{thm:01.1.3} 与 \ref{prop:01.1.4} 的证明，记 $\varphi_{jk}$ 为 $(\D_X,i_X)|_{U_j\cap U_k}$ 的自同构，使 $\psi_j=\varphi_{jk}\circ\psi_k$，并记
$$s_{jk}=\alpha_j(\alpha_k^{-1}(1))\in\calO_X^*,\qquad 1\le j<k\le n.$$
对 $1\le i\le n$，
$$\sigma_i=\alpha_i\otimes\psi_i:(L\otimes_{\calO_X}\D)|_{U_i}\longrightarrow\D_X|_{U_i}$$
是 $\calO_X$-模的同构，且对 $s\in L|_{U_i}$、$u,v\in\D|_{U_i}$ 有
$$\sigma_i(s\otimes uv)=\alpha_i(s)\psi_i(uv)=\alpha_i(s)\psi_i(u)\psi_i(v)=\sigma_i(s\otimes u)\psi_i(v).$$
因此，若通过映射 $T\mapsto T(1)$ 把「看作右 $\D_X$-模的 $\D_X$」的微分自同态与 $\D_X$ 等同起来，则得到自然同构 $\tau_i:\D_L|_{U_i}\to\D_X|_{U_i}$，由
$$\tau_i(u)=(\sigma_i\circ u\circ\sigma_i^{-1})(1)=\sigma_i\bigl(u(\alpha_i^{-1}(1)\otimes 1)\bigr)$$
给出。又对 $1\le j<k\le n$ 及任意 $s\in L|_{U_j\cap U_k}$、$v\in\D|_{U_j\cap U_k}$ 有
$$\sigma_j(s\otimes v)=\alpha_j(s)\psi_j(v)=s_{jk}\alpha_k(s)\varphi_{jk}(\psi_k(v))=s_{jk}\varphi_{jk}(\sigma_k(s\otimes v)),$$
这意味着对 $u\in\D_L|_{U_j\cap U_k}$ 有
\begin{align*}
\tau_j(u)&=\sigma_j\bigl(u(\alpha_j^{-1}(1)\otimes 1)\bigr)=s_{jk}\varphi_{jk}\bigl(\sigma_k(u(s_{jk}^{-1}\alpha_k^{-1}(1)\otimes 1))\bigr)\\
&=s_{jk}\varphi_{jk}\bigl(\sigma_k(u(\alpha_k^{-1}(1)\otimes s_{jk}^{-1}))\bigr)=s_{jk}\varphi_{jk}\bigl(\sigma_k(u(\alpha_k^{-1}(1)\otimes 1)s_{jk}^{-1})\bigr)\\
&=s_{jk}\varphi_{jk}\bigl(\sigma_k(u(\alpha_k^{-1}(1)\otimes 1))\,s_{jk}^{-1}\bigr)=s_{jk}\varphi_{jk}(\tau_k(u))\,s_{jk}^{-1}.
\end{align*}
若置
$$\rho_{jk}(v)=s_{jk}\varphi_{jk}(v)s_{jk}^{-1},\qquad v\in\D_X|_{U_j\cap U_k},$$
则得到 $\D_X|_{U_j\cap U_k}$ 的自同构 $\rho_{jk}$，使得 $\tau_j=\rho_{jk}\circ\tau_k$。如前，记 $\omega=(\omega_{jk};1\le j<k\le n)$ 为 $C^1(\calU,Z^1_X)$ 中对应于 $(\varphi_{jk};1\le j<k\le n)$ 的元素。设 $\xi\in T_X|_{U_j\cap U_k}$，则
$$\rho_{jk}(\xi)=s_{jk}\varphi_{jk}(\xi)s_{jk}^{-1}=s_{jk}(\xi-\omega_{jk}(\xi))s_{jk}^{-1}=\xi-s_{jk}^{-1}ds_{jk}(\xi)-\omega_{jk}(\xi),$$
因此 $C^1(\calU,Z^1_X)$ 中对应于 $(\rho_{jk};1\le j<k\le n)$ 的元素等于
$$(\omega_{jk}+s_{jk}^{-1}ds_{jk};\ 1\le j<k\le n).$$
\end{proof}

\begin{proposition}\label{prop:01.1.9}
设 $\D$ 是 $X$ 上的扭微分算子层，$L$ 与 $L'$ 是两个可逆 $\calO_X$-模。则扭微分算子层 $(\D_L)_{L'}$ 自然同构于 $\D_{L'\otimes_{\calO_X}L}$。
\end{proposition}

\begin{proof}
显然
$$L'\otimes_{\calO_X}L\otimes_{\calO_X}\D=(L'\otimes_{\calO_X}\D_L)\otimes_{\D_L}(L\otimes_{\calO_X}\D),$$
作为 $\calO_X$-模与右 $\D$-模。因此 $(\D_L)_{L'}$ 在第二式第一个因子上的右作用与 $\D$ 的右作用交换。这给出从 $(\D_L)_{L'}$ 到 $\D_{L'\otimes_{\calO_X}L}$ 的自然态射。由 \ref{lem:01.1.1}，这个态射是扭微分算子层的同构。
\end{proof}

以下我们通过这个同构把 $(\D_L)_{L'}$ 与 $\D_{L'\otimes_{\calO_X}L}$ 等同起来。

\begin{proposition}\label{prop:01.1.10}
设 $\phi:X\to Y$ 是光滑代数簇的态射，$\D$ 是 $Y$ 上的扭微分算子层，$L$ 是可逆 $\calO_Y$-模。则扭微分算子层 $(\D_L)_\phi$ 自然同构于 $(\D_\phi)_{\phi^*(L)}$。
\end{proposition}

\begin{proof}
显然
\begin{align*}
\D_{X\to Y}&=\phi^*(\D)=\phi^*(L^{-1}\otimes_{\calO_Y}L\otimes_{\calO_Y}\D)=\phi^*(L)^{-1}\otimes_{\calO_X}\phi^*(L\otimes_{\calO_Y}\D)\\
&=\phi^*(L)^{-1}\otimes_{\calO_X}(\D_L)_{X\to Y}\otimes_{\phi^{-1}(\D_L)}\phi^{-1}(L\otimes_{\calO_Y}\D)\\
&=(\phi^*(L)^{-1}\otimes_{\calO_X}(\D_L)_\phi)\otimes_{(\D_L)_\phi}(\D_L)_{X\to Y}\otimes_{\phi^{-1}(\D_L)}\phi^{-1}(L\otimes_{\calO_Y}\D),
\end{align*}
作为 $\calO_X$-模与右 $\phi^{-1}(\D)$-模。显然，$((\D_L)_\phi)_{\phi^*(L)}$ 在第一个因子上的左作用与 $\phi^{-1}(\D)$ 的右作用交换。因此存在从 $((\D_L)_\phi)_{\phi^*(L)}$ 到 $\D_\phi$ 的自然态射。由 \ref{lem:01.1.1}，它是扭微分算子层的同构。把这个自然同构用 $\phi^*(L)$ 去扭，并利用 \ref{prop:01.1.9}，即得：$(\D_L)_\phi$ 自然同构于 $(\D_\phi)_{\phi^*(L)}$。
\end{proof}

若 $A$ 是 $X$ 上的 $k$-代数层，记 $A^\circ$ 为 $X$ 上与之反同构的 $k$-代数层（\emph{反代数\index{fandaishu@反代数|hyperpage}层}）。

\begin{proposition}\label{prop:01.1.11}
设 $(\D,i)$ 是光滑代数簇 $X$ 上的扭微分算子层。则 $(\D^\circ,i)$ 也是 $X$ 上的扭微分算子层。
\end{proposition}

设 $n=\dim X$，$\omega_X$ 是 $X$ 上 $n$ 次微分形式构成的层。则李代数层 $T_X$ 在层 $\omega_X$ 上有自然作用：向量场 $\xi$ 通过相应的 Lie 导数 $L_\xi$ 作用。事实上，若 $U$ 是 $X$ 的开集，$\omega\in\omega_X(U)$，$\eta_1,\eta_2,\dots,\eta_n$ 是 $U$ 上的局部向量场，则
$$(L_\xi\omega)(\eta_1\wedge\eta_2\wedge\dots\wedge\eta_n)=\xi(\omega(\eta_1\wedge\dots\wedge\eta_n))-\sum_{i=1}^n\omega(\eta_1\wedge\dots\wedge[\xi,\eta_i]\wedge\dots\wedge\eta_n).$$
设 $\xi,\xi'$ 是 $U$ 上的局部向量场，则
{\small
\begin{align*}
([L_\xi,L_{\xi'}]\omega)(\eta_1\wedge\dots\wedge\eta_n)&=\xi((L_{\xi'}\omega)(\eta_1\wedge\dots\wedge\eta_n))-\xi'((L_\xi\omega)(\eta_1\wedge\dots\wedge\eta_n))\\
&\quad-\sum_{i=1}^n(L_{\xi'}\omega)(\eta_1\wedge\dots\wedge[\xi,\eta_i]\wedge\dots\wedge\eta_n)\\
&\quad+\sum_{i=1}^n(L_\xi\omega)(\eta_1\wedge\dots\wedge[\xi',\eta_i]\wedge\dots\wedge\eta_n)\\
&=[\xi,\xi'](\omega(\eta_1\wedge\dots\wedge\eta_n))+\sum_{i=1}^n\omega(\eta_1\wedge\dots\wedge[\xi',[\xi,\eta_i]]\wedge\dots\wedge\eta_n)\\
&\quad-\sum_{i=1}^n\omega(\eta_1\wedge\dots\wedge[\xi,[\xi',\eta_i]]\wedge\dots\wedge\eta_n)\\
&=[\xi,\xi'](\omega(\eta_1\wedge\dots\wedge\eta_n))\\
&\quad-\sum_{i=1}^n\omega(\eta_1\wedge\dots\wedge[[\xi,\xi'],\eta_i]\wedge\dots\wedge\eta_n)\\
&=(L_{[\xi,\xi']}\omega)(\eta_1\wedge\dots\wedge\eta_n),
\end{align*}
}
这说明上述作用是李代数层的作用。

又对 $U$ 上的正则函数 $f$ 与向量场 $\xi$ 有
\begin{align*}
L_\xi(f\omega)(\eta_1\wedge\dots\wedge\eta_n)&=\xi(f\omega(\eta_1\wedge\dots\wedge\eta_n))-\sum_{i=1}^n f\omega(\eta_1\wedge\dots\wedge[\xi,\eta_i]\wedge\dots\wedge\eta_n)\\
&=(\xi(f)\omega+fL_\xi(\omega))(\eta_1\wedge\dots\wedge\eta_n).
\end{align*}
由此立得 $[L_\xi,f]=\xi(f)$，即 $L_\xi$ 是 $\omega_X$ 上的一阶微分算子。

取足够小的 $U$，可设 $T_U$ 是自由 $\calO_U$-模，即存在 $U$ 上的局部向量场 $\eta_1,\eta_2,\dots,\eta_n$ 构成 $T_U$ 的 $\calO_U$-基。于是可把 $\xi$ 写成 $\xi=\sum_{i=1}^n g_i\eta_i$，其中 $g_i\in\calO_U$。这意味着
\begin{align*}
(L_{f\xi}\omega)(\eta_1\wedge\dots\wedge\eta_n)&=f\xi(\omega(\eta_1\wedge\dots\wedge\eta_n))-\sum_{i=1}^n\omega(\eta_1\wedge\dots\wedge[f\xi,\eta_i]\wedge\dots\wedge\eta_n)\\
&=f\xi(\omega(\eta_1\wedge\dots\wedge\eta_n))-\sum_{i=1}^n f\omega(\eta_1\wedge\dots\wedge[\xi,\eta_i]\wedge\dots\wedge\eta_n)\\
&\quad+\sum_{i=1}^n\omega(\eta_1\wedge\dots\wedge\underset{i}{\xi}\wedge\dots\wedge\eta_n)\eta_i(f)\\
&=\Bigl(fL_\xi(\omega)+\sum_{i=1}^n g_i\eta_i(f)\,\omega\Bigr)(\eta_1\wedge\dots\wedge\eta_n)\\
&=(fL_\xi(\omega)+\xi(f)\omega)(\eta_1\wedge\dots\wedge\eta_n)=L_\xi(f\omega)(\eta_1\wedge\dots\wedge\eta_n),
\end{align*}
即 $L_{f\xi}\,\omega=L_\xi(f\omega)$，亦即在 $\omega_X$ 上有 $L_{f\xi}=L_\xi{}f$（右乘 $f$）。
\begin{annote}
底本 p.9 这一行公式在文本层有明显串行脱漏（末式印作「$(fL_\xi(\omega)+\xi(f)\omega)(\eta_1\wedge\dots\wedge\eta_n)L_\xi(f\omega)(\eta_1\wedge\dots\wedge\eta_n)$」）。按推导补全为 $fL_\xi(\omega)+\xi(f)\omega=L_\xi(f\omega)$，这正是映射 $\lambda=-L$ 关于右乘的 $\calO_X$-线性。[来源：底本 p.9，按推导校正]
\end{annote}
这意味着映射 $\lambda:\xi\mapsto-L_\xi$ 是 $\calO_X$-模的态射，从 $T_X$ 映入「看作关于右乘的 $\calO_X$-模」的 $\D_{\omega_X}$。它具有性质
$$\lambda([\xi,\eta])=[\lambda(\eta),\lambda(\xi)],\qquad \xi,\eta\in T_X.$$
因此由 [9]，它可延拓为 $k$-代数层 $\D_X$ 到 $\D_{\omega_X}^\circ$ 的态射，且在 $\calO_X$ 上是恒等的。由 \ref{lem:01.1.1}，这说明它实际上是同构。

于是有下述结果。

\begin{lemma}\label{lem:01.1.12}
设 $X$ 是维数为 $n$ 的光滑代数簇，$\omega_X$ 是 $X$ 上 $n$ 次微分形式构成的可逆 $\calO_X$-模。则二元组 $(\D_X^\circ,i_X)$ 自然同构于 $(\D_{\omega_X},i_{\omega_X})$。
\end{lemma}

由这个结果立即得到 \ref{prop:01.1.11}。因此可以计算 $\D^\circ$ 的同构类。

\begin{proposition}\label{prop:01.1.13}
设 $X$ 是光滑代数簇，$\D$ 是 $X$ 上的扭微分算子层。则
$$t(\D^\circ)=-t(\D)+H^1(d\log)(i(\omega_X)).$$
\end{proposition}

\begin{proof}
取 $X$ 的开覆盖 $\calU=(U_i;1\le i\le n)$ 及相应的同构 $\psi_i:(\D,i)|_{U_i}\to(\D_{U_i},i_{U_i})$。如前，对 $1\le j<k\le n$ 记 $\varphi_{jk}$ 为 $(\D_{U_j\cap U_k},i_{U_j\cap U_k})$ 的自同构，使 $\psi_j=\varphi_{jk}\circ\psi_k$。设 $\omega_{jk}$ 是由 \ref{lem:01.1.2} 由 $\varphi_{jk}$ 决定的闭 $1$-形式，即对 $U_j\cap U_k$ 上任意局部向量场 $\xi$ 有 $\varphi_{jk}(\xi)=\xi-\omega_{jk}(\xi)$。则 $\psi_i$（$1\le i\le n$）也是 $(\D^\circ,i)|_{U_i}$ 到 $(\D_X^\circ,i_X)|_{U_i}$ 的同构。与 \ref{prop:01.1.11} 证明中引入的映射 $\lambda$ 复合，得到同构
$$\tau_i:(\D^\circ,i)|_{U_i}\longrightarrow(\D_{\omega_X},i_{\omega_X})|_{U_i},\qquad 1\le i\le n.$$
此外，自同构 $\varphi_{jk}$ 通过 $\sigma_{jk}=\lambda\circ\varphi_{jk}\circ\lambda^{-1}$ 定义了 $(\D_{\omega_X},i_{\omega_X})|_{U_j\cap U_k}$ 的自同构，使得 $\tau_j=\sigma_{jk}\circ\tau_k$（$1\le j<k\le n$）。显然，$\sigma_{jk}$ 由
$$\sigma_{jk}(\lambda(\xi))=\lambda(\varphi_{jk}(\xi))=\lambda(\xi-\omega_{jk}(\xi))=\lambda(\xi)-\omega_{jk}(\xi)$$
确定，其中 $\xi$ 是 $U_j\cap U_k$ 上的任意局部向量场。不妨设开集 $U_i$ 足够小，使得存在 $\calO_{U_i}$-模同构 $\alpha_i:\omega_{U_i}\to\calO_{U_i}$（$1\le i\le n$）。它们通过
$$\beta_i(\eta)=\alpha_i\circ\eta\circ\alpha_i^{-1},\qquad \eta\in\D_{\omega_{U_i}},$$
定义同构 $\beta_i:(\D_{\omega_{U_i}},i_{\omega_{U_i}})\to(\D_{U_i},i_{U_i})$（$1\le i\le n$）。又如 \ref{prop:01.1.4} 的证明，置
$$s_{jk}=\alpha_j(\alpha_k^{-1}(1)),\qquad 1\le j<k\le n.$$
复合 $\gamma_i=\beta_i\circ\tau_i$ 是 $(\D^\circ,i)|_{U_i}$ 到 $(\D_{U_i},i_{U_i})$ 的同构（$1\le i\le n$）。自同构
$$\delta_{jk}=\beta_j\circ\sigma_{jk}\circ\beta_k^{-1},\qquad 1\le j<k\le n,$$
满足
$$\gamma_j=\beta_j\circ\tau_j=\beta_j\circ\sigma_{jk}\circ\tau_k=\delta_{jk}\circ\gamma_k.$$
设 $\xi$ 是 $U_j\cap U_k$ 上的局部向量场。则对 $f\in\calO_X$ 有
\begin{align*}
\beta_i(\lambda(\xi))(f)&=\alpha_i(\lambda(\xi)\alpha_i^{-1}(f))=\alpha_i(\lambda(\xi)(f\alpha_i^{-1}(1)))\\
&=-\alpha_i(L_\xi(f\alpha_i^{-1}(1)))=-\xi(f)+\beta_i(\lambda(\xi))(1),
\end{align*}
因此
$$\beta_i^{-1}(\xi)=-\lambda(\xi)+\beta_i(\lambda(\xi))(1).$$
这给出
\begin{align*}
\delta_{jk}(\xi)&=\beta_j\bigl(\sigma_{jk}(\beta_k^{-1}(\xi))\bigr)=-\beta_j(\sigma_{jk}(\lambda(\xi)))+\beta_k(\lambda(\xi))(1)\\
&=-\beta_j(\lambda(\varphi_{jk}(\xi)))+\beta_k(\lambda(\xi))(1)=-\beta_j(\lambda(\xi))+\omega_{jk}(\xi)+\beta_k(\lambda(\xi))(1)\\
&=\xi+\omega_{jk}(\xi)-\beta_j(\lambda(\xi))(1)+\beta_k(\lambda(\xi))(1).
\end{align*}
如同 \ref{prop:01.1.4} 的证明，可见 $\beta_j(\lambda(\xi))=s_{jk}\beta_k(\lambda(\xi))s_{jk}^{-1}$，这意味着
$$\beta_j(\lambda(\xi))(1)=s_{jk}\beta_k(\lambda(\xi))(s_{jk}^{-1})=-s_{jk}\xi(s_{jk}^{-1})+\beta_k(\lambda(\xi))(1),$$
最终得到
\begin{align*}
\delta_{jk}(\xi)&=\xi+\omega_{jk}(\xi)-\beta_j(\lambda(\xi))(1)+\beta_k(\lambda(\xi))(1)\\
&=\xi+\omega_{jk}(\xi)-s_{jk}^{-1}ds_{jk}(\xi)=\xi-\bigl(-\omega_{jk}+s_{jk}^{-1}ds_{jk}\bigr)(\xi).
\end{align*}
\end{proof}

\section{齐性扭微分算子层\index{qixingniuweifensuanziceng@齐性扭微分算子层|hyperpage}}
\label{sec:01.1.2}

设 $G$ 是特征 $0$ 的代数闭域 $k$ 上的连通代数群，$X$ 是它的齐性空间\index{qixingkongjian@齐性空间|hyperpage}。对 $G$ 在 $X$ 的结构层 $\calO_X$ 上的作用求微分，得到代数同态
$$\tau:U(\g)\longrightarrow\Gamma(X,\D_X).$$
显然，这个映射是 $G$-等变的。

设 $\D$ 是 $X$ 上的扭微分算子层，带有 $G$ 的代数作用 $\gamma$ 以及代数同态 $\alpha:U(\g)\to\Gamma(X,\D)$，使得

\begin{enumerate}[label=(\roman*)]
  \item $\D$ 中的乘法是 $G$-等变的；
  \item $G$ 在 $\D$ 上作用的微分与作用 $T\mapsto[\alpha(\xi),T]$（$\xi\in\g$，$T\in\D$）一致；
  \item 映射 $\alpha:U(\g)\to\Gamma(X,\D)$ 是 $G$-模的态射。
\end{enumerate}

则称三元组 $(\D,\gamma,\alpha)$ 是 $X$ 上的\emph{齐性扭微分算子层}。本节将分类 $X$ 上所有的齐性扭微分算子层。

显然，$\D_X$ 连同 $G$ 的自然作用与同态 $\tau$ 给出 $X$ 上的齐性扭微分算子层。

在向量空间层 $U^\circ=\calO_X\otimes_k U(\g)$ 上，可以通过置
$$\xi(f\otimes\eta)=\tau(\xi)f\otimes\eta+f\otimes\xi\eta,\qquad \xi\in\g,\ \eta\in U(\g),\ f\in\calO_X,$$
来定义 $U(\g)$-模张量积的结构。另一方面，$U^\circ=\calO_X\otimes_k U(\g)$ 通过在第一个因子上的乘法具有 $\calO_X$-模的结构。并且对 $\xi\in\g$、$\eta\in U(\g)$、$f,g\in\calO_X$ 有
$$[\xi,g](f\otimes\eta)=\xi(gf\otimes\eta)-g\,\xi(f\otimes\eta)=\tau(\xi)(g)f\otimes\eta=[\tau(\xi),g]f\otimes\eta.$$
这意味着 $U(\g)$ 通过微分算子作用在 $U^\circ$ 上，而相应的同态 $\Psi:U(\g)\to\Diff(U^\circ)$（$\Diff(U^\circ)$ 是 $U^\circ$ 上微分算子构成的环）与按次数的滤过相容。可把 $\Psi$ 延拓为 $\calO_X$-模的态射 $\calO_X\otimes_k U(\g)\to\Diff(U^\circ)$，它把 $f\otimes\xi$（$f\in\calO_X$，$\xi\in U(\g)$）对应于微分算子 $f\Psi(\xi)\in\Diff(U^\circ)$。对 $f\in\calO_X$、$\xi\in U(\g)$ 有
$$\Psi(f\otimes\xi)(1\otimes 1)=f\Psi(\xi)(1\otimes 1)=f\otimes\xi,$$
这意味着 $\Psi:\calO_X\otimes_k U(\g)\to\Diff(U^\circ)$ 是单射。我们断言它的像是 $\Diff(U^\circ)$ 的子环层。显然，关于左乘它是 $\calO_X$-模，关于右乘它是右 $\Psi(U(\g))$-模。因此只需证明：对任意 $f\in\calO_X$、$\xi\in\g$，微分算子 $\Psi(\xi)f$ 也在 $\Psi$ 的像中。而
$$\Psi(\xi)f=[\Psi(\xi),f]+f\Psi(\xi)=\tau(\xi)f+f\Psi(\xi),$$
末式显然属于 $\Psi(\calO_X\otimes_k U(\g))$。这意味着 $U^\circ$ 具有自然的环层结构，使得 $\Psi:U^\circ\to\Diff(U^\circ)$ 是同态。此外，乘法由
$$(f\otimes\xi)(g\otimes\eta)=f\tau(\xi)g\otimes\eta+fg\otimes\xi\eta$$
给出，其中 $f,g\in\calO_X$，$\xi\in\g$，$\eta\in U(\g)$。由此可见 $\tau$ 可延拓为环层 $U^\circ$ 到 $\D_X$ 的同态。令 $\g^\circ=\calO_X\otimes_k\g$，看作 $U^\circ$ 的 $\calO_X$-子模。则对任意 $f,g\in\calO_X$、$\xi,\eta\in\g$ 有
$$[f\otimes\xi,g\otimes\eta]=f\tau(\xi)g\otimes\eta-g\tau(\eta)f\otimes\xi+fg\otimes[\xi,\eta],$$
这意味着 $\g^\circ$ 在此运算下是李代数层。由这个计算可见，$\tau$ 定义了 $\g^\circ$ 到 $X$ 上局部向量场层 $T_X$ 的同态，我们仍用 $\tau$ 记之。

\begin{lemma}\label{lem:01.2.1}
态射 $\tau:\g^\circ\to T_X$ 是满态射\index{mantaishe@满态射|hyperpage}。
\end{lemma}

\begin{proof}
$\calO_X$-模 $\g^\circ$ 与 $T_X$ 都是局部自由的，故结论归结为：态射 $\tau$ 在任意 $x\in X$ 处诱导的几何纤维\index{jihexianwei@几何纤维|hyperpage}上的线性映射 $T_x(\tau)$ 是满射。
\end{proof}

可以通过置
$$F_pU^\circ=\calO_X\otimes_k F_pU(\g),\qquad p\in\Z_+,$$
在 $U^\circ$ 上定义递增滤过，其中 $FU(\g)$ 是包络代数 $U(\g)$ 的标准滤过。显然，这个滤过与 $U^\circ$ 上的代数结构以及同态 $\tau:U^\circ\to\D_X$ 都相容。并且
$$F_0U^\circ=\calO_X,\qquad F_1U^\circ=\calO_X\oplus\g^\circ,$$
且 $F_1U^\circ$ 作为代数层生成 $U^\circ$。

记 $\mathfrak{b}^\circ$ 为 $\tau:\g^\circ\to T_X$ 的核。则 $\mathfrak{b}^\circ$ 是 $\g^\circ$ 中的理想层\index{lixiangceng@理想层|hyperpage}。此外，若 $\sum_i f_i\otimes\xi_i\in\mathfrak{b}^\circ$，$g\in\calO_X$，则
$$\Bigl[\sum_i f_i\otimes\xi_i,\ g\otimes 1\Bigr]=\sum_i f_i\tau(\xi_i)g\otimes 1=0;$$
这意味着 $J_0=\mathfrak{b}^\circ U^\circ$ 是 $U^\circ$ 中的双边理想层。

\begin{proposition}\label{prop:01.2.2}
\begin{enumerate}[label=(\roman*)]
  \item 态射 $\tau:U^\circ\to\D_X$ 是满态射。
  \item 态射 $\tau:U^\circ\to\D_X$ 的核是理想层 $J_0$。
\end{enumerate}
\end{proposition}

\begin{proof}
(i) 由 \ref{lem:01.1.1} 以及 $\D_X$ 由 $\calO_X$ 与 $T_X$ 生成得出。

(ii) 显然 $J_0=\mathfrak{b}^\circ U^\circ$ 含于 $\tau$ 的核中。又对任意 $x\in X$，几何纤维 $T_x(J_0)=\mathfrak{b}_xU(\g)$ 正是线性映射 $T_x(\tau)$ 的核，这里 $T_x(\tau)$ 把 $U^\circ$ 的几何纤维 $T_x(U^\circ)=U(\g)$ 映入 $\D_X$ 在 $x$ 处的几何纤维。
\end{proof}

现在对 $X$ 上的齐性扭微分算子层证明类似的结果。

设 $(\D,\gamma,\alpha)$ 是 $X$ 上的齐性扭微分算子层。则由 (ii)，对任意 $\xi\in\g$、$f\in\calO_X$ 有
$$[\alpha(\xi),f]=[\tau(\xi),f]=\tau(\xi)f.$$
特别地，对任意 $f,g\in\calO_X$ 有 $[[\alpha(\xi),f],g]=0$，故 $\alpha(\xi)$ 的次数不超过 $1$（$\xi\in\g$）。定义映射 $U^\circ\to\D$ 为 $f\otimes T\mapsto f\alpha(T)$（$f\in\calO_X$，$T\in U(\g)$），并沿用记号 $\alpha$ 记之。则对任意 $f,g\in\calO_X$、$\xi,\eta\in\g$ 有
\begin{align*}
\alpha((f\otimes\xi)(g\otimes\eta))&=\alpha(f\tau(\xi)g\otimes\eta+fg\otimes\xi\eta)=f\tau(\xi)(g)\,\alpha(\eta)+fg\,\alpha(\xi\eta)\\
&=f[\alpha(\xi),g]\alpha(\eta)+fg\,\alpha(\xi)\alpha(\eta)=f\alpha(\xi)g\alpha(\eta)=\alpha(f\otimes\xi)\alpha(g\otimes\eta).
\end{align*}
因此 $\alpha$ 延拓为环层的态射，且与自然滤过相容。由于 $\Gr\D=S(T_X)=\Gr\D_X$，且它作为 $\calO_X$-代数由 $\Gr_1\D$ 生成，我们立即得到 $\Gr\alpha=\Gr\tau$，并且 $\Gr\alpha:\Gr U^\circ\to\Gr\D$ 是环层的满态射。这意味着 $\alpha(\mathfrak{b}^\circ)\subset\calO_X$，即 $\alpha$ 定义了 $G$-齐性 $\calO_X$-模 $\mathfrak{b}^\circ$ 到 $\calO_X$ 的 $G$-等变态射 $\sigma$。

固定基点\index{jidian@基点|hyperpage} $x_0\in X$。其稳定化子\index{wendinghuazi@稳定化子|hyperpage} $B_0$ 作用在 $\mathfrak{b}_0$ 的对偶空间 $\mathfrak{b}_0^*$ 上。记 $I(\mathfrak{b}_0^*)$ 为 $\mathfrak{b}_0^*$ 中 $B_0$-不变元构成的子空间。则 $I(\mathfrak{b}_0^*)$ 与「$G$-齐性 $\calO_X$-模 $\mathfrak{b}^\circ$ 到 $\calO_X$ 的全部 $G$-等变态射 $\sigma$ 构成的空间」之间有自然的线性同构。因此 $(\D,\gamma,\alpha)$ 决定 $I(\mathfrak{b}_0^*)$ 中的一个元素。

对每个 $\lambda\in I(\mathfrak{b}_0^*)$ 可联系一个 $G$-等变态射 $\sigma_\lambda:\mathfrak{b}^\circ\to\calO_X$。设 $\phi_\lambda:\mathfrak{b}^\circ\to U^\circ$ 由
$$\phi_\lambda(s)=s-\sigma_\lambda(s),\qquad s\in\mathfrak{b}^\circ$$
给出。则 $\im\phi_\lambda$ 生成 $U^\circ$ 中的双边理想层 $J_\lambda$。置
$$\D_{X,\lambda}=U^\circ/J_\lambda.$$
这是 $X$ 上的代数层。

\begin{proposition}\label{prop:01.2.3}
代数层 $\D_{X,\lambda}$ 是 $X$ 上的扭微分算子层。
\end{proposition}

称 $\D_{X,\lambda}$ 为 $X$ 上\emph{联系于 $\lambda$ 的齐性扭微分算子层}。

作为前述讨论与 \ref{prop:01.2.2} 的推论，有下述结果。

\begin{theorem}\label{thm:01.2.4}
映射 $\lambda\mapsto\D_{X,\lambda}$ 是 $I(\mathfrak{b}_0^*)$ 到 $X$ 上齐性扭微分算子层的同构类集合上的双射。
\end{theorem}

\begin{proof}
设 $(\D,\gamma,\alpha)$ 是 $X$ 上的齐性扭微分算子层。则由前述讨论，它决定唯一的 $\lambda\in I(\mathfrak{b}_0^*)$。此外，$J_\lambda$ 含于同态 $\alpha:U^\circ\to\D$ 的核中。这意味着 $\alpha$ 诱导环层的同态 $\beta:\D_{X,\lambda}\to\D$，它与 $\D_{X,\lambda}$ 和 $\D$ 的滤过相容，并与 $U(\g)$ 到 $\Gamma(X,\D_{X,\lambda})$ 和 $\Gamma(X,\D)$ 的自然映射相容。又 $\Gr\beta$ 是分次环层的同构，由此 $\beta$ 也是同构。
\end{proof}

为证 \ref{prop:01.2.3}，由 \ref{prop:01.2.2} 与 $G$-齐性，只需找到基点 $x_0$ 的一个邻域 $U$ 以及 $U^\circ|_U$ 的局部自同构 $\Psi_\lambda$，使得 $\Psi_\lambda|_{\calO_U}=1$ 且 $\Psi_\lambda(J_0|_U)=J_\lambda|_U$。

设 $U$ 是 $X$ 的开集。我们要描述 $U^\circ|_U$ 的一些具有下述性质的自同构 $\rho$：

\begin{enumerate}[label=(\roman*)]
  \item 对任意 $f\in\calO_U$ 有 $\rho(f)=f$；
  \item $\Gr\rho$ 是恒等。
\end{enumerate}

显然，$\rho$ 由它在 $1\otimes\xi$（$\xi\in\g$）上的值完全确定。又 (ii) 意味着
$$\rho(1\otimes\xi)=1\otimes\xi-\omega(\xi)\otimes 1,\qquad \omega(\xi)\in\calO(U).$$
由 (i) 还得
$$\rho(f\otimes\xi)=f\otimes\xi-f\omega(\xi)\otimes 1,\qquad f\in\calO_U,\ \xi\in\g.$$
要成为自同构，$\rho$ 还必须满足
$$\rho([1\otimes\xi,1\otimes\eta])=[\rho(1\otimes\xi),\rho(1\otimes\eta)]=1\otimes[\xi,\eta]-\tau(\xi)\omega(\eta)\otimes 1+\tau(\eta)\omega(\xi)\otimes 1,$$
即
\begin{equation}\label{eq:01.2.1}
\omega([\xi,\eta])=\tau(\xi)\omega(\eta)-\tau(\eta)\omega(\xi)
\end{equation}
对任意 $\xi,\eta\in\g$ 成立。因此 $\omega$ 是从 $\g$ 到 $\calO(U)$ 的线性映射，且被以 $\calO(U)$ 为系数的 $\g$ 的李代数上同调的微分零化。此外，可把 $\omega$ 延拓为 $\calO_U$-模的态射 $\g^\circ|_U\to\calO_U$，由 $\omega(f\otimes\xi)=f\omega(\xi)$（$f\in\calO_U$，$\xi\in\g$）给出。关系 \eqref{eq:01.2.1} 意味着
\begin{align*}
\omega([f\otimes\xi,g\otimes\eta])&=fg\omega([\xi,\eta])+f\tau(\xi)(g)\omega(\eta)-g\tau(\eta)(f)\omega(\xi)\\
&=\tau(f\otimes\xi)(\omega(g\otimes\eta))-\tau(g\otimes\eta)(\omega(f\otimes\xi)),
\end{align*}
即对任意两个截面 $s,s'\in\g^\circ|_U$ 有
\begin{equation}\label{eq:01.2.2}
\omega([s,s'])=\tau(s)(\omega(s'))-\tau(s')(\omega(s)).
\end{equation}
还要指出 $\omega$ 是局部的，即若 $s\in\g^\circ$ 使 $s(x)=0$（$x\in U$），则 $\omega(s)(x)=0$。此外，由 \eqref{eq:01.2.2}，对 $s\in\mathfrak{b}^\circ$、$s'\in\g^\circ|_U$ 有
$$\omega([s',s])=\tau(s')(\omega(s)),$$
这意味着映射 $\omega:\mathfrak{b}^\circ|_U\to\calO_U$ 是 $\g^\circ|_U$-模的态射。

下面这个结果将要用到。

\begin{lemma}\label{lem:01.2.5}
设 $\zeta:\mathfrak{b}^\circ|_U\to\calO_U$ 是局部 $\g^\circ|_U$-模态射，$x\in U$。若对任意 $s\in\mathfrak{b}^\circ$ 有 $\zeta(s)(x)=0$，则存在 $x$ 的邻域 $V\subset U$，使 $\zeta|_V=0$。
\end{lemma}

\begin{proof}
设 $\xi\in\g$。则
$$\bigl(\tau(\xi)(\zeta(s))\bigr)(x)=\zeta([1\otimes\xi,s])(x)=0.$$
由此可见 $\zeta(s)$ 在 $x$ 处的各阶导数都为零，故 $\zeta(s)$ 在 $x$ 处的芽为零。由 $\mathfrak{b}^\circ$ 的凝聚\index{ningju@凝聚|hyperpage}性可知 $\zeta$ 在 $x$ 的某个邻域上为零。
\end{proof}

另一方面，若线性映射 $\omega:\g\to\calO(U)$ 满足关系 \eqref{eq:01.2.1}，则它通过
$$\rho_\omega(f\otimes\xi)=f\otimes\xi-f\omega(\xi)\otimes 1,\qquad f\in\calO_U,\ \xi\in\g,$$
定义了 $U^\circ|_U$ 的自同构 $\rho_\omega$，且满足 (i) 与 (ii)。显然，所有这样的 $\omega$ 构成一个向量空间。

下面构造一些满足上述性质的映射 $\omega$。

(I) 设 $\chi$ 是 $B_0$ 的特征标\index{tezhengbiao@特征标|hyperpage}。设 $s$ 是齐性可逆 $\calO_X$-模 $\calO(\chi)$ 在 $U$ 上的截面。对 $\xi\in\g$ 置
$$\xi s=\omega(\xi)s.$$
则
\begin{align*}
\omega([\xi,\eta])s&=[\xi,\eta]s=\xi(\eta s)-\eta(\xi s)=\xi(\omega(\eta)s)-\eta(\omega(\xi)s)\\
&=\tau(\xi)(\omega(\eta))s-\tau(\eta)(\omega(\xi))s+\omega(\eta)\xi s-\omega(\xi)\eta s\\
&=\bigl(\tau(\xi)(\omega(\eta))-\tau(\eta)(\omega(\xi))\bigr)s,
\end{align*}
即 $\omega$ 满足所需条件。

(II) 设
$$0\longrightarrow k\longrightarrow V\longrightarrow k\longrightarrow 0$$
是 $B_0$ 的代数表示的正合列，其中 $B_0$ 在 $k$ 上平凡地作用。设
$$0\longrightarrow\calO_X\longrightarrow\mathcal{V}\longrightarrow\calO_X\longrightarrow 0$$
是相应的 $G$-齐性局部自由 $\calO_X$-模的正合列。设 $s'$ 是 $\mathcal{V}$ 中作为 $\calO_X$ 的截面 $1$ 的像的截面，$s''$ 是 $\mathcal{V}$ 的局部截面，使它在 $x_0$ 处的芽映入 $1$ 的芽。则存在 $x_0$ 的邻域 $U$，使得 $s'$ 与 $s''$ 构成 $\mathcal{V}|_U$ 作为 $\calO_U$-模的一组基，且对任意 $\xi\in\g$ 有
$$\xi s''=\omega(\xi)s'\qquad\text{和}\qquad \xi s'=0.$$
于是
$$\omega([\xi,\eta])s'=[\xi,\eta]s''=\xi(\omega(\eta)s')-\eta(\omega(\xi)s')=\bigl(\tau(\xi)(\omega(\eta))-\tau(\eta)(\omega(\xi))\bigr)s'.$$
因此 $\omega$ 再次具有所需性质。

现在来证 \ref{prop:01.2.3}。首先，$\mathfrak{b}_0$ 上任一 $B_0$-不变的线性形式在 $[\mathfrak{b}_0,\mathfrak{b}_0]$ 上为零。设 $C_0$ 是 $B_0$ 的换位子群的单位分支。则 $C_0$ 是 $B_0$ 的闭正规子群。商群 $D_0=B_0/C_0$ 是单位分支交换的代数群，其李代数为 $\mathfrak{d}_0=\mathfrak{b}_0/[\mathfrak{b}_0,\mathfrak{b}_0]$。设 $D_0=L_0U_0$ 是 $D_0$ 的 Levi 分解。则 $U_0$ 是交换的幂幺\index{miyao@幂幺|hyperpage}子群，而 $L_0$ 的单位分支是 $D_0$ 的单位分支中全部半单\index{bandan@半单|hyperpage}元构成的环面\index{huanmian@环面|hyperpage}。因此有直和分解 $\mathfrak{d}_0=\mathfrak{l}_0\oplus\mathfrak{u}_0$，且两个直和项都是 $D_0$-不变的。这意味着由 $\lambda$ 定义的 $\mathfrak{d}_0$ 上 $D_0$-不变线性形式 $\mu$ 可以写成两个 $D_0$-不变线性形式 $\mu_1$ 与 $\mu_2$ 之和，它们分别在 $\mathfrak{l}_0$ 与 $\mathfrak{u}_0$ 上为零。把这些线性形式与 $\mathfrak{b}_0$ 到 $\mathfrak{d}_0$ 的投影复合，即得到 $\lambda$ 分解为 $\lambda_1$ 与 $\lambda_2$ 之和。可以定义 $U_0$ 在 $k^2$ 上的表示，使 $u\in U_0$ 通过矩阵
$$\begin{pmatrix}1&\mu_1(\log u)\\ 0&1\end{pmatrix}$$
作用；显然它可延拓为 $D_0$ 的表示，其中 $L_0$ 平凡地作用。此外，可把它解释为 $B_0$ 的表示。应用 (II) 中的构造，可在 $x_0$ 的邻域 $U$ 上构造满足 \eqref{eq:01.2.1} 的线性映射 $\omega_1:\g\to\calO(U)$，使得 $\omega_1|_{\mathfrak{b}_0}=\lambda_1$。

另一方面，$\mathfrak{l}_0$ 上任一线性形式都是 $L_0$ 的单位分支的特征标的微分的线性组合。通过用 $L_0$ 的连通分支群取平均，可知 $\mathfrak{l}_0$ 上任一 $L_0$-不变线性形式都是 $L_0$-不变特征标（即 $L_0$ 的一维表示的微分）的线性组合。应用 (I) 中的构造，可在 $x_0$ 的邻域 $U$ 上得到满足 \eqref{eq:01.2.1} 的线性映射 $\omega_2:\g\to\calO(U)$，使得 $\omega_2|_{\mathfrak{b}_0}=\lambda_2$。

于是，在 $x_0$ 的邻域 $U$ 上得到满足 \eqref{eq:01.2.1} 的线性映射 $\omega:\g\to\calO(U)$，使得 $\omega|_{\mathfrak{b}_0}=\lambda$。相应的 $\g^\circ|_U$-模态射 $\omega:\mathfrak{b}^\circ|_U\to\calO_U$ 由 \ref{lem:01.2.5} 在 $x_0$ 的某个更小的邻域 $V$ 上与 $\sigma_\lambda$ 一致。这又意味着 $\rho_\omega$ 是 $U^\circ|_V$ 的自同构，满足 $\rho_\omega|_{\calO_V}=1$ 且 $\rho_\omega(J_0|_V)=J_\lambda|_V$。
